\documentclass[11pt]{article}
\usepackage[margin=1in]{geometry}
\usepackage{amsmath,amssymb,amsthm,hyperref}
\newtheorem{theorem}{Theorem}
\newtheorem{lemma}[theorem]{Lemma}
\begin{document}
\setlength{\emergencystretch}{2em}
\title{Lower-two binomial SFC at prime gaps}
\author{}\date{}\maketitle
The mixed-rank identity below holds at every gap. At prime gaps it
proves all-index lower-two binomial SFC for every positive real
algebraic Gamma shape integral at one place above that prime.
The argument imports the monic factorized operator, adjoint
solutions and full concomitant rank from the
\href{HIGHER_GAP_ADJOINT_RANK.tex}{higher-gap adjoint note}, equations (1)--(8) and the proof
of \texttt{thm:rank}, especially its full-jet argument.
No lower-root simplicity or uniform rank theorem is assumed.

Write $\mathcal L_\beta(X^j)=(\beta)_j$ for the normalized
Gamma moment functional. Fix $d\ge2$, $n\ge1$, $\gamma>0$, and set
\[
 T_{n,\gamma,d}(t)=\sum_{k=0}^n\binom nk(\gamma)_{dk}t^k,\quad
 \epsilon=(-1)^{d+1},\quad
 R=z^{dn}T_{n,\gamma,d}(\epsilon z^{-d}),\quad
 M=\gamma+dn,\quad E=M-z\partial_z.
\]
Products $(E)_j$ are rising products with step one.
Write $B_j=B_{j,\gamma,d}$, always at the SAME $\gamma$, with
\[
 B_0=1,\qquad B_{j+1}=(z+dj+\gamma+d-1)B_j-zB_j'.
\]
The cyclic column of $J$ at $u$ is its remainder coefficient
column modulo $z^d-u$, in basis $1,z,\ldots,z^{d-1}$.

\begin{theorem}[Offset-two mixed rank at every gap]\label{thm:mixed}
Put
\[
 A=(M+d)(M+d+1),\qquad \kappa=\frac{d(n+1)}A,\qquad
 H=B_{n+d-1}+\kappa z^dB_n.
\]
For $u\ne0$, the two equalities
\[
 T_{n,\gamma,d}(\epsilon/u)
   =T_{n+1,\gamma+2,d}(\epsilon/u)=0
\]
hold if and only if the $d$ by $(d-1)$ cyclic matrix
\[
 [B_{n+1},B_{n+2},\ldots,B_{n+d-2},H](u)
\]
has rank less than $d-1$. When $d=2$ the consecutive list is empty.
After removal of the nonzero gauges the identity is algebraic
in characteristic zero.
\end{theorem}
\begin{proof}
The lower Euler equation is
\[
 \theta(E)_dR=\epsilon z^d(dn-\theta)R
             =\epsilon z^d(E-\gamma)R,\qquad \theta=z\partial_z.
\]
Pascal's identity and Gamma shifting give
\[
 (\gamma)_2z^{d(n+1)}
       T_{n+1,\gamma+2,d}(\epsilon z^{-d})
   =[z^d(E)_2+\epsilon(E)_{d+2}]R.
\]
The commutator $Ez^d=z^d(E-d)$ reduces its right hand side to
\[
 \epsilon A(E)_dR+
 z^d[\gamma(M+d+1)-d(n+1)E]R.
\]
Indeed $E(E)_dR=M(E)_dR-\epsilon z^d(E-\gamma)R$;
applying $E$ once more proves the reduction. At $R=0$, the upper
moment vanishes exactly when
\[
 (E)_dR=\epsilon\kappa u\,ER.                         \tag{1}
\]

Here are the concomitant columns explicitly.
The monic lower operator is
\[
 L=z^{-d}\partial_z z^{M+d}\partial_z^d z^{-M}
       -\partial_z+dn/z.
\]
For $f_\xi=e^{-\xi z}z^{1-M}B_n(\xi z)$, $\xi^d=1$, put
$S=z^{-M}R$ and $h_\xi=z^{M+d}(z^{-d}f_\xi)'$.
Repeated integration by parts gives the Lagrange concomitant
\[
 {\cal C}_\xi=
 f_\xi z^MS^{(d)}
 -\sum_{j=0}^{d-1}(-1)^jh_\xi^{(j)}S^{(d-1-j)}
 -f_\xi R.                                           \tag{2}
\]
Differentiating (2) gives $f_\xi LR-RL^*f_\xi$.
The imported adjoint identity and exponential-times-rational
argument make it identically zero for this $R$.
The full concomitant matrix has rank $d$ on the $(d+1)$-jet,
with kernel the line spanned by the nonzero jet of $R$,
as in the imported full-jet proof.

The $B$ recurrence gives the exact simplification
\[
 h_\xi=-e^{-\xi z}B_{n+1}(\xi z).
\]
Since
$\partial_z^k(z^{-M}R)=(-1)^kz^{-M-k}(E)_kR$,
on $R=0$, division of (2) by the nonzero gauge
$(-1)^de^{-\xi z}z^{1-M-d}$ gives
\[
 B_n(w)(E)_dR-
   \sum_{j=0}^{d-2}L_j(w)(E)_{d-1-j}R=0,\qquad w=\xi z,       \tag{3}
\]
where
\[
 L_j(w)=w^j(\partial_w-1)^jB_{n+1}(w).
\]
These are the full rising-$E$ jet columns.
They triangularize by CONSTANT column operations, because
\[
 L_{j+1}=w(\partial_w-1)L_j-jL_j,
\]
\[
 w(\partial_w-1)B_{n+1+r}
    =(M+2d-1+dr)B_{n+1+r}-B_{n+2+r}.
\]
Inductively $L_j$ is a constant linear combination of
$B_{n+1},\ldots,B_{n+j+1}$, with last coefficient $(-1)^j$.
This induction holds for every $j$.

Impose (1) in (3). The $ER$ column is
$-L_{d-2}+\epsilon\kappa u B_n$; modulo the other columns it is
\[
 (-1)^{d-1}B_{n+d-1}+\epsilon\kappa u B_n=\epsilon H.
\]
The remaining columns triangularize to the stated consecutive list.
The restrictions $R=0$ and (1) are independent and define a
$(d-1)$-dimensional jet subspace. Full concomitant rank makes a
nonzero restricted kernel equivalent to the ACTUAL $R$ jet satisfying
both restrictions, proving both directions without simplicity.
Finally Fourier evaluation at $\xi z$ is invertible in characteristic
zero when $z\ne0$, and yields the cyclic matrix.
No reduction of that Fourier matrix modulo a prime is used.
\end{proof}

\begin{lemma}[Affine residual polynomial]\label{lem:univalent}
Let $K$ be a complete algebraically closed valued field. If
$F\in\mathcal O_K[x]$ reduces modulo its maximal ideal to $c(1+x)$
with $c\ne0$, then $F$ is injective on $\mathcal O_K$.
\end{lemma}
\begin{proof}
For integral $x,y$, the polynomial difference quotient
$(F(x)-F(y))/(x-y)$ reduces to $c$, hence is a unit.
This uses a finite polynomial identity, not a convergent series.
\end{proof}

\begin{theorem}[Lower exponent two at every prime gap]\label{thm:prime}
Let $p$ be prime and let $\beta>0$ be real algebraic.
Suppose some embedding
$\iota:\mathbb Q(\beta)\longrightarrow\mathbb C_p$ satisfies
$v_p(\iota(\beta))\ge0$. At every $n\ge1$ the two polynomials
\[
 T_{n,2n+\beta,p},\qquad T_{n+1,2n+\beta+2,p}
\]
are coprime. In particular this holds for every positive rational
$p$-integral $\beta$, including $\beta=1$.
Thus $X^2(A+BX^p)$ has a nonzero Gamma-$\beta$ moment in every
pair of consecutive positive orders whenever $(A,B)\ne(0,0)$.
\end{theorem}
\begin{proof}
Use $\iota$ and suppress it from the notation.
Put $\gamma=2n+\beta$ and $\epsilon=(-1)^{p+1}$.
The reversed polynomial $u^nT_{n,\gamma,p}(\epsilon/u)$ is monic
with integral coefficients. A root $t$ of $T$ therefore gives
$u=\epsilon/t$ integral. Constant terms are one, so $t,u\ne0$.

For any integral $u$, the reduced cyclic space is
$k[z]/(z^p-\bar u)$. The auxiliary recurrence reduces to
\[
 {\mathsf A}=z-z\partial_z+\bar\gamma-1.
\]
It has cyclic vector $1$, since the first $p$ powers are monic of
successive degrees. Its cyclic bidiagonal matrix has characteristic
polynomial
\[
 \det(xI-{\mathsf A})=
 \prod_{j=0}^{p-1}(x-\bar\gamma+1+j)-\bar u
 =x^p-x-Q,\quad Q=\bar u+\bar\gamma^p-\bar\gamma.
\]
Hence $\mathsf A^p=\mathsf A+Q$.
If $Q\ne0$, $\mathsf A$ is invertible, so every $p-1$
consecutive powers of its cyclic vector are independent.
If $Q=0$, its eigenvalues are the distinct elements of $\mathbb F_p$,
and the cyclic vector has a nonzero component in every eigenspace.
The $p-1$ nonzero eigenspaces and a Vandermonde determinant show
that every $p-1$ consecutive POSITIVE powers are independent here too.

First suppose $\bar\gamma\notin\mathbb F_p$.
The denominator factors $(M+p),(M+p+1)$ reduce to
$\bar\gamma,\bar\gamma+1$, both nonzero, so $\kappa$ has positive
valuation. Since $u$ is integral, the final mixed column reduces
to $\mathsf A^{n+p-1}1$ without perturbation.
All $p-1$ columns are consecutive positive powers, so the mixed
matrix has full column rank. No common root is possible.

It remains that $\bar\gamma\in\mathbb F_p$.
Choose an integer $c$ with $v_p(\gamma-c)>0$ and set
$\delta=\min\{1,v_p(\gamma-c)\}>0$; use $\delta=1$ if $\gamma=c$.
In each block of $p$ consecutive factors of $(\gamma)_{pk}$,
one has valuation at least $\delta$ and the others are integral.
The coefficient at degree $k$ in $T_{n,\gamma,p}$ has valuation
at least $k\delta$. The unique least-term argument for the reversed
monic polynomial gives $v_p(u)\ge\delta>0$.
Now $Q=0$ and $\mathsf A^p=\mathsf A$.

Put $\eta=\kappa u$. If $v_p(\eta)>0$, the mixed reduction consists
of the independent powers
$\mathsf A^{n+1}1,\ldots,\mathsf A^{n+p-1}1$.
If $v_p(\eta)<0$, scale the last column by $\eta^{-1}$.
The reduction is
$\mathsf A^{n+1}1,\ldots,\mathsf A^{n+p-2}1,\mathsf A^{n}1$,
again independent. Thus a common root forces $v_p(\eta)=0$
and $q=\kappa^{-1}$ to have positive valuation.

Choose a row set with nonzero maximal minor of this latter
Krylov matrix. Let $F(\eta)$ be the corresponding EXACT minor
of the cyclic coefficient columns
\[
 [B_{n+1}(q\eta),\ldots,B_{n+p-2}(q\eta),
                   B_{n+p-1}(q\eta)+\eta B_n(q\eta)].
\]
Its coefficients are integral. Since $q$ reduces to zero and
$\mathsf A^{n+p-1}1=\mathsf A^{n}1$ for $n\ge1$,
\[
 \bar F(\eta)=c_0(1+\eta),\qquad c_0\ne0.
\]
Lemma~\ref{lem:univalent} gives at most one integral zero of $F$.
Every actual common root, and EVERY conjugate of it over
$\mathbb Q(\beta)$, satisfies the same identities and integrality
balance. Distinct roots give distinct $\eta=\epsilon\kappa/t$.
In characteristic zero all conjugates must therefore coincide,
forcing the common root into $\mathbb Q(\beta)$.

Its original real embedding makes it real, contradicting positivity
of the even-index member at every real $t$:
\[
 T_{j,2j+\beta,p}(t)=\frac1{\Gamma(2j+\beta)}
       \int_0^\infty x^{2j+\beta-1}(1+tx^p)^j e^{-x}\,dx>0.
\]
The inequality is asserted for even $j$; one member has even index.
The phase normalization is
\[
 {\cal L}_\beta((X^2(A+BX^p))^j)
       =A^j(\beta)_{2j}T_{j,2j+\beta,p}(B/A).
\]
When a coefficient vanishes the nonzero phase is a monomial and all
positive moments are nonzero.

For $p=2$ the intermediate list is EMPTY: the matrix is the single
two-component column $B_{n+1}+\kappa uB_n$.
A nonzero component of $\mathsf A^n1$ supplies the scalar minor.
The proof remains valid with $u=-1/t$ throughout.
Fourier inversion is used only in characteristic zero before
cyclic coefficient reduction.
\end{proof}
No unrestricted all-binomial or arbitrary-support SFC assertion is made.
No algebraic shape of negative valuation at every $p$-adic embedding
is covered. The retained gap-two theorem has a stronger
all-positive-real-shape scope.
\end{document}
